% ID: FIS-TER-CAL-002
% Titolo: Riscaldamento dell'acqua con pannelli solari
% Disciplina: Fisica
% Area: Termologia
% Argomento: Calore specifico
% Sottoargomento: Riscaldamento di una massa d'acqua
% Classe: Seconda liceo scientifico
% Difficolta: 3
% Tipo: problema_numerico
% Risultato: soluzione da aggiungere
% Tempo_stimato: 8 min
% Competenze: bilancio energetico, calore specifico, conversione del tempo
% Prerequisiti: calore specifico, energia, potenza
% Tag: termologia, calore_specifico, acqua, energia_termica, pannelli_solari
% Fonte: verifica_fisica_termologia_anonimizzata
% Autore: Riccardo Carli
% Licenza: CC-BY-SA-4.0

\begin{esercizio}
I pannelli solari di un impianto ricevono dal Sole
\(1{,}8 \cdot 10^4\,\mathrm{J}\) di calore ogni secondo e usano l'80\% di
questo calore per riscaldare \(495\,\mathrm{kg}\) di acqua inizialmente a
\(12\,^\circ\mathrm{C}\).

Calcola la temperatura dell'acqua dopo 1 ora di funzionamento.
\end{esercizio}

\begin{soluzione}
% Soluzione da aggiungere.
\end{soluzione}

